\documentclass[notes=show]{beamer}
%%%%%%%%%%%%%%%%%%%%%%%%%%%%%%%%%%%%%%%%%%%%%%%%%%%%%%%%%%%%%%%%%%%%%%%%%%%%%%%%%%%%%%%%%%%%%%%%%%%%%%%%%%%%%%%%%%%%%%%%%%%%%%%%%%%%%%%%%%%%%%%%%%%%%%%%%%%%%%%%%%%%%%%%%%%%%%%%%%%%%%%%%%%%%%%%%%%%%%%%%%%%%%%%%%%%%%%%%%%%%%%%%%%%%%%%%%%%%%%%%%%%%%%%%%%%
\AtBeginSubsection[
  {\frame<beamer>{\frametitle{Outline}   
    \tableofcontents[currentsection,currentsubsection]}}%
]%
{
  \frame<beamer>{ 
    \frametitle{Outline}   
    \tableofcontents[currentsection,currentsubsection]}
}
\usepackage{ulem}
\usepackage{pdflscape}
\setbeamertemplate{caption}{\insertcaption}
\usepackage{color,xcolor,ucs}
\usepackage{beamerthemesplit}
\usepackage{graphicx}
\usepackage{amsmath}
\usepackage{adjustbox}
\usepackage{pdfpages} 
\definecolor{mintgreen}{rgb}{0.6, 1.0, 0.6}
\makeatletter
\@addtoreset{subfigure}{framenumber}% subfigure counter resets every frame
\makeatother
\newcommand*{\Scale}[2][4]{\scalebox{#1}{$#2$}}%
\newcommand*{\Resize}[2]{\resizebox{#1}{!}{$#2$}}%
\usepackage{amsfonts}
\usepackage{booktabs}
\usepackage{color,soul}
\usepackage{pdfpages}
\setboolean{@twoside}{false}
\usepackage{amsmath}
\usepackage{pdflscape} 
\usepackage{eurosym}
\usepackage{amstext} % for \text
\DeclareRobustCommand{\officialeuro}{%
  \ifmmode\expandafter\text\fi
  {\fontencoding{U}\fontfamily{eurosym}\selectfont e}}
\usepackage[caption = false]{subfig}
\usepackage{appendixnumberbeamer} 
\usepackage{graphicx}
\usepackage{mathpazo}
\usepackage{hyperref}
\usepackage{multimedia}
\usepackage{graphicx}
\usepackage{multirow}
\usepackage{subfig}
\usepackage{verbatim}
\usepackage{booktabs, multicol, multirow}
\setcounter{MaxMatrixCols}{10}
\input{Macros.tex}
\AtBeginSection[]
{
  \begin{frame}[noframenumbering]
    \frametitle{Outline for section \thesection}
    \tableofcontents[currentsection]
  \end{frame}
}

\usepackage[style=british]{csquotes}
\usepackage{xcolor,soul}
\renewcommand<>{\hl}[1]{\only#2{\beameroriginal{\hl}}{#1}}

% https://tex.stackexchange.com/questions/41683/why-is-it-that-coloring-in-soul-in-beamer-is-not-visible
\makeatletter
\newcommand\SoulColor{%
  \let\set@color\beamerorig@set@color
  \let\reset@color\beamerorig@reset@color}
\makeatother
\SoulColor
\AtBeginSection{\frame{\sectionpage}}

\newsavebox\mybox
\newenvironment{aquote}[1]
  {\savebox\mybox{#1}\begin{quote}\openautoquote\hspace*{-.7ex}}
  {\unskip\closeautoquote\vspace*{1mm}\signed{\usebox\mybox}\end{quote}}

\newtheorem{proposition}[theorem]{Proposition}
\newtheorem{Assumption}[theorem]{Assumption}
\newenvironment{stepenumerate}{\begin{enumerate}[<+->]}{\end{enumerate}}
\newenvironment{stepitemize}{\begin{itemize}[<+->]}{\end{itemize} }
\newenvironment{stepenumeratewithalert}{\begin{enumerate}[<+-| alert@+>]}{\end{enumerate}}
\newenvironment{stepitemizewithalert}{\begin{itemize}[<+-| alert@+>]}{\end{itemize} }
\usetheme{Madrid}
\makeatletter
\setbeamertemplate{footline}
{
  \leavevmode%
  \hbox{%
  \begin{beamercolorbox}[wd=.333333\paperwidth,ht=2.25ex,dp=1ex,center]{author in head/foot}%
    \usebeamerfont{author in head/foot}\insertshortauthor~~\beamer@ifempty{\insertshortinstitute}{}{(\insertshortinstitute)}
  \end{beamercolorbox}%
  \begin{beamercolorbox}[wd=.333333\paperwidth,ht=2.25ex,dp=1ex,center]{title in head/foot}%
    \usebeamerfont{title in head/foot}\insertshorttitle
  \end{beamercolorbox}%
  \begin{beamercolorbox}[wd=.333333\paperwidth,ht=2.25ex,dp=1ex,right]{date in head/foot}%
    \usebeamerfont{date in head/foot}\insertshortdate{}\hspace*{2em}
%    \insertframenumber{} / \inserttotalframenumber\hspace*{2ex} % DELETED
  \end{beamercolorbox}}%
  \vskip0pt%
}
\makeatother
\usepackage{pdfpages}
\usepackage[style=british]{csquotes}
\title{Firms as Wage Setters}
\author{Raffaele Saggio}
\date{UBC}

\def\signed #1{{\leavevmode\unskip\nobreak\hfil\penalty50\hskip1em
  \hbox{}\nobreak\hfill #1%
  \parfillskip=0pt \finalhyphendemerits=0 \endgraf}}

% (removed duplicate \mybox / aquote definition)
%\AtBeginSection[]
%{
%  \begin{frame}
%    \frametitle{Table of Contents}
%    \tableofcontents[currentsection]
%  \end{frame}
%}

\providecommand{\stretchon}{}\providecommand{\stretchoff}{} % no-op: package not installed
\begin{document}
\begin{frame}[plain]
	\titlepage
\end{frame}

\stretchon
\begin{frame}
\frametitle{Introduction}
\begin{itemize}
\item What drives variation in wages?
\bigskip
\item Traditional view: 
\medskip
\begin{enumerate}
\pause
\item Trends in wage inequality driven by changes in market level skill prices. \\ {\tiny{Katz and Autor (1999); Acemoglu and Autor, (2011)}}.
\pause
\medskip
\item Labor market institutions might provide barriers to price adjustments. \\ {\tiny{DiNardo, Fortin and Lemieux (1994).}}
\pause
\medskip
\item No role for firms.
\end{enumerate}
\bigskip
\pause
\item ``Monopsony" view: 
\begin{enumerate}
\bigskip
\item Employers have considerable latitude to set wages. {\tiny{Robison 1933}}. 
\medskip
\item Firm heterogeneity influences not only who is hired but also the corresponding level of pay.
\end{enumerate}
\pause
\bigskip
\item Today: discuss the importance of firms in wage settings.
\medskip
\begin{itemize}
\item Much progress in recent years (matched employer-employee datasets).
\end{itemize}
\end{itemize}
\end{frame}
%%
\begin{frame}
\frametitle{Yet they move (together)}
\begin{columns}
	\column{\dimexpr\paperwidth+5pt}
	\begin{center}
		\includegraphics[scale=0.35]{aux/barth}
	\end{center}
\end{columns}
\end{frame}
%%

\begin{frame}[shrink=20]
\frametitle{Outline}
\begin{enumerate}
\item Empirical studies of Industry Wage Differentials.
\medskip
\begin{itemize}
\item Slitcher (1950);
\medskip
\item Krueger and Summers (1989); 
\medskip
\item Gibbons and Katz (1992).
\medskip
\end{itemize}
\pause
\medskip
\item The arrival of matched employer-employee datasets: 
\pause
\medskip
\begin{itemize}
\item Abowd, Kramarz and Margolis (1999)
\medskip
\item Card, Heining and Kline (2013).
\medskip
\item Song, Price, Guvenen, Bloom and von Wachter (2018).
\end{itemize}
\medskip
\pause
\item Microfounding the role of Firms in the Labor Market.
\begin{itemize}
\medskip
\item Card, Cardoso, Heining and Kline (2018).
\medskip
\item Sorkin (2018).
\end{itemize}
\medskip
\item Applications.
\begin{itemize}
\item Gender gaps.
\item Temp/Outsourced workers.
\item Job displacement.
\end{itemize}
\end{enumerate}
\end{frame}
\begin{frame}\centering
		\includegraphics[scale=0.5]{aux/joke}
\end{frame}
\begin{frame}[plain]
\frametitle{Where It All Started? Slitcher (1950)}
\begin{columns}
	\column{\dimexpr\paperwidth+5pt}
	\begin{center}
		\includegraphics[scale=0.3]{aux/Slichter}
	\end{center}
\end{columns}
\end{frame}
%%
\section{Industry Wage Differentials}
%%
\begin{frame}[plain]
\frametitle{Slitcher (1950)}
\begin{itemize}
\medskip
\item Studies the wage structure of low skilled men in Boston using plant-level survey. 
\begin{columns}
	\column{\dimexpr\paperwidth+5pt}
	\begin{center}
		\includegraphics[scale=0.28]{aux/Slichter2}
	\end{center}
\end{columns}

\end{itemize}
\end{frame}
\begin{frame}[plain]
\frametitle{Findings}
\begin{itemize}
\item Slitcher (1950) documents 7 empirical regularities for wages of unskilled men:
\bigskip
\begin{enumerate}
\item Positive correlation with wages of skilled co-workers.
\item Negative correlation with \% females.
\item Positive correlation with industry value added / worker hour.
\item Positive correlation with sales / worker hour.
\item Negative correlation with payroll / sales.
\item Positive correlation with value added / sales.
\item High correlation of industry wage rank over time.
\end{enumerate}
\bigskip
\item Conclusion: managerial wage policies are important to understand dispersion of wages across industries.
\pause
\bigskip
\bigskip
\item Discussion:
\medskip
\begin{itemize}
\item Worker pay correlates with productivity differences.
\medskip
\item Does this mean that firms share ``rents" with workers? 
\medskip
\item or these correlations simply reflect worker sorting?
\end{itemize}
\end{itemize}
\hyperlink{ks}{\beamerbutton{Industry Wage Differentials literature}} 
\end{frame}
%%%
\begin{frame}
\label{ks}
\frametitle{Krueger and Summers (1989)}
\pause
\begin{itemize}
\medskip
\item Exploit longitudinal variation on workers.
\medskip
\item ``Experiment": workers switching industries $\rightarrow$ industry wage premium 
\pause
\medskip
\item Question: are OLS estimates of industry wage premia $\approx$ to Fixed effects estimates?
\medskip
\item Useful to assess biases due to sorting in measuring industry wage premia.
\end{itemize}
\end{frame}
%%
\begin{frame}
\frametitle{Very similar estimates across methods!}
\begin{columns}
	\column{\dimexpr\paperwidth+5pt}
	\begin{center}
		\includegraphics[scale=0.45]{aux/krueger}
	\end{center}
\end{columns}
\end{frame}
%%
\begin{frame}
\frametitle{No evidence of compensating differentials}
\begin{columns}
	\column{\dimexpr\paperwidth+5pt}
	\begin{center}
		\includegraphics[scale=0.22]{aux/krueger2}
	\end{center}
\end{columns}
\end{frame}
%%
\begin{frame}
\frametitle{People don't quit high paying industries}
\begin{columns}
	\column{\dimexpr\paperwidth+5pt}
	\begin{center}
		\includegraphics[scale=0.29]{aux/krueger3}
	\end{center}
\end{columns}
\end{frame}
%%
\begin{frame}
\frametitle{Gibbons and Katz (1992): Endogenous mobility}
\begin{itemize}
\item Fixed Effects estimates from ``switchers" design can still be biased 
\medskip
\item Idea:
\medskip
\begin{itemize}
\medskip
\item Talent is unknown ex-ante. 
\medskip
\item When talent is revealed, the good worker moves to the high productive industry and gets a wage bump.
\medskip
\item Interpretation: not rents getting shared but revelation of talent + sorting on match effects!
\end{itemize}
\medskip
\pause
\item Solution: Exploit ``exogenous" separations. 
\medskip
\item Plant closings in CPS Displaced Worker Survey.
\medskip
\item Same exercise as Krueger and Summers (1989): compare FE with OLS.
\end{itemize}
\end{frame}
%%

%%
\begin{frame}
\frametitle{FE estimates $\approx$ Cross Sectional Estimates}
\begin{columns}
	\column{\dimexpr\paperwidth+5pt}
	\begin{center}
		\includegraphics[scale=0.275]{aux/katz}
	\end{center}
\end{columns}
\end{frame}
\begin{frame}
\frametitle{Summing up}
\begin{itemize}
\item (Plant Closing sample): Slope is 0.79 with $R^{2}=0.72$.
\medskip
\item (Layoff sample): Slope is 0.97 with $R^{2}=0.81$.
\medskip
\item Evidence is therefore: 
\medskip
\begin{itemize}
\item Hard to reconcile with a story of talent revelation/sorting on match effects.
\item More in line with a true ``industry wage effects model". 
\end{itemize}
\end{itemize}
\end{frame}
%%
%%
\section{Evidence from Matched Employer-Employee Data}
%%
%%%
\begin{frame}
\frametitle{Abowd, Kramarz and Margolis (1999)}
\begin{itemize}
\item Questions:
\begin{enumerate}
\medskip
\item Are there big pay differences across firms within industry?
\item Does where you work ultimately matter for how much you earn?
\medskip
\end{enumerate}
\item The AKM model (Two-Way Fixed Effects Model)
\be
y_{it}=\underbrace{\alpha_{i}}_{\text{Person Effects}}+\underbrace{\psi_{J(i,t)}}_{\text{Firm Effects}}+ W_{it}^{\top}\gamma+r_{it}
\ee
\item $\alpha_{i}\equiv$ worker ``ability", rewarded equally at all firms.
\item $\psi_{J(i,t)}\equiv$ Firm effect, i.e. firm $j$ pays constant $\psi_{j}$ above/below reference firm to \textit{all employees}.  
\item $W_{it}^{\top}\gamma\equiv$ economy wide changes in returns to education/experience.  
\pause
\medskip
\item Immense computational challenge for 90's technologies. 
\end{itemize}
\end{frame}
%%
\begin{frame}
\frametitle{Two-Way Fixed Effects}
\begin{itemize}
\item Rewrite model in matrix notation
\be
\underbrace{y}_{n\times 1} = \underbrace{D}_{n \times G}\alpha +  \underbrace{F}_{n \times J} \psi + \underbrace{W}_{n \times P}\gamma + \underbrace{r}_{n\times 1}  = X\beta + r
\ee
\item
\medskip $T_{g}\equiv$ Total number of periods in which we observe worker $g$.
\item $n=\sum_{g=1}^{G}T_{g}$ $\equiv$ Total number of observations.
\medskip
\item $K=G + J + P$ $\equiv$ Total number of parameters to be estimated. 
\medskip
\item AKM framework extends to various setups (doctors-patients, teachers-students, etc).
\medskip
\item Builds on two fundamental assumptions:
\be
\label{KEY}
\underbrace{\E[\varepsilon \vert D, F] = 0}_{\text{Exogenous Mobility}}; \qquad \underbrace{rank(X^{\top}X)=K-1}_{\text{Invertibility}}; 
\ee 
\end{itemize}
\end{frame}
%%
\begin{frame}
\frametitle{Abowd, Creecy and Kramarz (2002): Invertibility Condition}
\begin{columns}
	\column{\dimexpr\paperwidth+5pt}
	\begin{center}
		\includegraphics[scale=0.3]{aux/akm0}
	\end{center}
\end{columns}
\pause
\begin{itemize}
\medskip
\item Can I identify $\psi_{3}-\psi_{1}$? (Suppose no-covariates for simplicity)
\pause
\medskip
\item Yes! Just use  $\E[\Delta y_{3}+ \Delta y_{1}]$
\end{itemize}
\end{frame}
%%
%%
\begin{frame}
\frametitle{Invertibility Condition}
\begin{columns}
	\column{\dimexpr\paperwidth+5pt}
	\begin{center}
		\includegraphics[scale=0.3]{aux/akm0}
	\end{center}
\end{columns}
\pause
\begin{itemize}
\medskip
\item Can I identify $\psi_{5}-\psi_{1}$?
\pause
\medskip
\item No!
\end{itemize}
\end{frame}
%%
%%
\begin{frame}
\frametitle{Focus estimation on the largest connected set}
\begin{columns}
	\column{\dimexpr\paperwidth+5pt}
	\begin{center}
		\includegraphics[scale=0.3]{aux/akm00}
	\end{center}
\end{columns}
\pause
\begin{itemize}
\item In the example above it is possible to identify all firm effects (after normalizing one firm effect to 0).
\pause
\medskip
\item \underline{Punch Line}: (High-Dimensional): Two-Way fixed effects models identified only on connected sets. 
\medskip
\item Several algorithms can find connected sets in enormous datasets in very short time
\end{itemize}
\end{frame}
%%
%%

\begin{frame}
\frametitle{Results from Abowd, Creecy and Kramarz (2002)}
\begin{columns}
	\column{\dimexpr\paperwidth+5pt}
	\begin{center}
		\includegraphics[scale=0.24]{aux/akm1}
	\end{center}
\end{columns}
\end{frame}
\begin{frame}
\frametitle{Firm effects: explain 30\% of the variance of earnings in France (!)}
\begin{columns}
	\column{\dimexpr\paperwidth+5pt}
	\begin{center}
		\includegraphics[scale=0.24]{aux/akm2}
	\end{center}
\end{columns}
\end{frame}
\begin{frame}
\frametitle{Correlation of Firm and Person effects: Negative (!)}
\begin{columns}
	\column{\dimexpr\paperwidth+5pt}
	\begin{center}
		\includegraphics[scale=0.24]{aux/akm3}
	\end{center}
\end{columns}
\end{frame}
\begin{frame}
\frametitle{Discussion}
\begin{itemize}
\item Abowd, Creecy and Kramarz (2002) find that firms wage policies explain b/w 20-30\% of the variation in wages.
\medskip
\item Worker and Firm fixed effects are negatively correlated $\rightarrow$ negative assortative matching?
\medskip
\item Abowd, Creecy and Kramarz (2002) studies a 1/25th sample for France. 1/10th sample from Washington data.
\medskip
\item Abowd, Lengermann and McKinney (2003): use 100\% file from LEHD (7 States).
\bigskip
\begin{enumerate}
\item Firm effects explain 16\% of the variance in wages.
\item Now firm and worker effects are positively correlated.
\end{enumerate}
\item Why this change in the estimates? Manifestation of sampling error/Limited Mobility biases $\rightarrow$ will get back to this in the applied econometrics part of the course. 

\end{itemize}
\end{frame}
\begin{frame}
\frametitle{Card, Heining and Kline (2013)}
\begin{itemize}
\item Universe of German social security records 1985-2009 (more than 400 million person-year observations)
\medskip
\item Studies \underline{changes} of the German wage structure. 
\medskip
\pause
\item This paper shifted the paradigm ($>$1000 GS citations):
\begin{itemize}
\item Before: Supply + Demand + Institutions.
\item CHK: Need S + D + I + F where F stands for Firms/Frictions. 
\end{itemize}
\medskip
\item Variance decomposition from AKM
\begin{enumerate}
\item Estimated over 4 different intervals.
\item  Each interval provides a ``snapshot" of the labor market.
\end{enumerate}
\medskip
\item Questions:
\begin{enumerate}
\item Did sorting change? 
\item Did importance of firms change?
\end{enumerate}
\bigskip
\item Discussion on endogenous mobility.
\end{itemize}
\end{frame}
%%
\begin{frame}
\frametitle{Wage Inequality in Germany}
\begin{columns}
	\column{\dimexpr\paperwidth+5pt}
	\begin{center}
		\includegraphics[scale=0.35]{aux/chk1}
	\end{center}
\end{columns}
\end{frame}
%%
%%
\begin{frame}
\frametitle{Discussion on Endogenous Mobility}
\begin{itemize}
\item AKM model
\be
y_{it} = \alpha_{i} + \psi_{J(i,t)} + W_{it}^{\top}\gamma + r_{it} 
\ee
\item CHK characterize the structural error of the AKM model as follows:
\be
r_{it}  = \underbrace{m_{i,J(i,t)}}_{\text{Match Effects}} + \underbrace{\nu_{it}}_{\text{Unit Root Component}} + \underbrace{\eta_{it}}_{\text{Transitory Component}}
\ee
\item AKM relies on exogenous mobility assumption- see equation (\ref{KEY})
\medskip
\item Sufficient condition:
\be
\label{ASS}
\Pr(J(i,t) = j \vert r) = \Pr(J(i,t) = j) = G_{jt}(\alpha_{i},\psi_{1}, \hdots \psi_{J}) \qquad \forall (i,t)
\ee
\pause
\item When is this assumption \underline{not} satisfied?
\end{itemize}
\end{frame}
%%
%%
\begin{frame}
\frametitle{Sources of Endogenous Mobility: Match Component}
\begin{itemize}
\item Sorting based on matched component. {\tiny{Roy (1951)}, Gibbons and Katz (1992)}
\begin{itemize}
\medskip
\item Wage changes for movers based on matched component:
\be
\begin{aligned}
& \E[y_{it}-y_{i,t-1}\vert J(i,t) = 2, J(i,t-1) =1]  = \psi_{2} - \psi_{1} +  \\ 
& + \underbrace{\E[m_{i2}-m_{i1} \vert J(i,t) = 2, J(i,t-1) =1 ]}_{>0 \text{  under positive Roy selection}}
\end{aligned}
\ee
\be
\begin{aligned}
& \E[y_{it}-y_{i,t-1}\vert J(i,t) = 1, J(i,t-1) =2]  = \psi_{1} - \psi_{2} +  \\ 
& + \underbrace{\E[m_{i1}-m_{i2} \vert J(i,t) = 1, J(i,t-1) =2 ]}_{>0 \text{  under positive Roy selection}}
\end{aligned}
\ee
\item AKM instead says
\be
\begin{aligned}
& \E[y_{it}-y_{i,t-1}\vert J(i,t) = 2, J(i,t-1) =1] = \\
& = - \E[y_{it}-y_{i,t-1}\vert J(i,t) = 1, J(i,t-1) =2]
\end{aligned}
\ee
\pause
\item We can test for that!
\end{itemize}
\end{itemize}
\end{frame}
%\appendix
%%
\begin{frame}
\frametitle{CHK Event Study }
\begin{columns}
	\column{\dimexpr\paperwidth+5pt}
	\begin{center}
		\includegraphics[scale=0.3]{aux/chk3}
	\end{center}
\end{columns}
\end{frame}
\begin{frame}
\frametitle{Wage changes when moving from 4-1}
\begin{columns}
	\column{\dimexpr\paperwidth+5pt}
	\begin{center}
		\includegraphics[scale=0.3]{aux/chk4}
	\end{center}
\end{columns}
\end{frame}
\begin{frame}
\frametitle{Wage changes when moving from 1-4: Roughly symmetric!}
\begin{columns}
	\column{\dimexpr\paperwidth+5pt}
	\begin{center}
		\includegraphics[scale=0.3]{aux/chk5}
	\end{center}
\end{columns}
\end{frame}
%%
\begin{frame}
\frametitle{Sources of Endogenous Mobility: Unit Root Component}
\begin{itemize}
\medskip
\item Drift in the expected wage individual $i$ can receive in period $t$ $(\nu_{it})$ predicts $\Pr(J(i,t) = j)$.  
\medskip
\item Example: Learning models (Gibbons et al., 2005).
\medskip
\item  Worker $i$ is revealed to be more productive than expected while employed by firm $j$:
\medskip
\begin{itemize}
\item Upward trend in wages while still employed at $j$.
\medskip
\item More likely to move in firms that employ high skilled workers. 
\end{itemize}
\pause
\medskip
\item We can test for that! 
\end{itemize}
\end{frame}
\begin{frame}
\frametitle{Back CHK Event Study }
\begin{columns}
	\column{\dimexpr\paperwidth+5pt}
	\begin{center}
		\includegraphics[scale=0.3]{aux/chk3}
	\end{center}
\end{columns}
\end{frame}
\begin{frame}
\frametitle{No Systematic trends prior to a move!}
\begin{columns}
	\column{\dimexpr\paperwidth+5pt}
	\begin{center}
		\includegraphics[scale=0.3]{aux/chk6}
	\end{center}
\end{columns}
\end{frame}
\begin{frame}
\frametitle{CHK decomposition of Germany Wage Inequality}
\begin{columns}
	\column{\dimexpr\paperwidth+5pt}
	\begin{center}
		\includegraphics[scale=0.28]{aux/chk7}
	\end{center}
\end{columns}
\end{frame}
\begin{frame}
\frametitle{CHK decomposition of Germany Wage Inequality}
\begin{columns}
	\column{\dimexpr\paperwidth+5pt}
	\begin{center}
		\includegraphics[scale=0.28]{aux/chk8}
	\end{center}
\end{columns}
\end{frame}
\begin{frame}
\frametitle{CHK decomposition of Germany Wage Inequality}
\begin{columns}
	\column{\dimexpr\paperwidth+5pt}
	\begin{center}
		\includegraphics[scale=0.28]{aux/chk9}
	\end{center}
\end{columns}
\end{frame}
\begin{frame}
\frametitle{CHK decomposition of Germany Wage Inequality}
\begin{columns}
	\column{\dimexpr\paperwidth+5pt}
	\begin{center}
		\includegraphics[scale=0.28]{aux/chk10}
	\end{center}
\end{columns}
\end{frame}
\begin{frame}
\frametitle{CHK decomposition of Germany Wage Inequality}
\begin{columns}
	\column{\dimexpr\paperwidth+5pt}
	\begin{center}
		\includegraphics[scale=0.28]{aux/chk11}
	\end{center}
\end{columns}
\end{frame}
\begin{frame}
\frametitle{CHK decomposition of Germany Wage Inequality}
\begin{columns}
	\column{\dimexpr\paperwidth+5pt}
	\begin{center}
		\includegraphics[scale=0.28]{aux/chk12}
	\end{center}
\end{columns}
\end{frame}
\begin{frame}
\frametitle{Jobs in High Firm Effects Firms Last Longer }
\begin{columns}
	\column{\dimexpr\paperwidth+5pt}
	\begin{center}
		\includegraphics[scale=0.35]{aux/chk13}
	\end{center}
\end{columns}
\end{frame}
\begin{frame}
\frametitle{Change in Returns to College Explained by Changes in Firm Effects}
\begin{columns}
	\column{\dimexpr\paperwidth+5pt}
	\begin{center}
		\includegraphics[scale=0.4]{aux/chk14}
	\end{center}
\end{columns}
\end{frame}
\begin{frame}
\frametitle{Cohort Effects in Variation of Firm Effects}
\begin{columns}
	\column{\dimexpr\paperwidth+5pt}
	\begin{center}
		\includegraphics[scale=0.4]{aux/chk15A}
	\end{center}
\end{columns}
\end{frame}
\begin{frame}
\frametitle{New firms are not applying sector-wide collective bargaining agreement?}
\begin{columns}
	\column{\dimexpr\paperwidth+5pt}
	\begin{center}
		\includegraphics[scale=0.4]{aux/chk15B}
	\end{center}
\end{columns}
\end{frame}
\begin{frame}
\frametitle{Distribution of Firm Effects across type of collective bargaining agreements}
\begin{columns}
	\column{\dimexpr\paperwidth+5pt}
	\begin{center}
		\includegraphics[scale=0.4]{aux/chk16}
	\end{center}
\end{columns}
\end{frame}
\begin{frame}
\frametitle{CHK Takeaway}
\begin{itemize}
\item AKM as an econometric tool to decompose the wage structure.
\medskip
\begin{itemize} 
\item Decompose gap in wages b/w groups into worker vs. firm component.
\medskip
\item Endogenous mobility not so bad?
\end{itemize}
\bigskip
\pause
\item Inserting the F in the S+D+I paradigm is crucial to understand wage inequality in Germany!
\medskip
\begin{itemize}
\item Firms are becoming more important (wage effects + sorting).
\medskip
\item Coincides with institutional changes.
\medskip
\item Cohort effects in the variation of firm effects. 
\end{itemize}
\end{itemize}
\end{frame}
%%
%%%%%
\section{Microfounding the Role of Firms in the Labor Market}
\begin{frame}
\frametitle{Introduction}
\begin{itemize}
\item We have seen that where you work matters for understanding wage inequality.
\bigskip
\item This echoes the original monopsony model of Robinson: a firm might have a firm-specific labor supply curve that permits the firm to set a wage that might be different from the ``market" wage.
\bigskip
\item But why firms might have a firm-specific labor supply curve? Key elements: (i) idiosyncratic taste heterogeneity (remember the cereal box examples) (ii) inability of firms to observe these idiosyncratic taste factors (no price discrimination).
\pause
\bigskip
\item We now present a ``modern" monospony model that incorporates (i) and (ii) and thus can ``microfound" the presence of a firm-specific labor supply function.   
\bigskip
\item Finally, we discuss whether differences in pay across firms simply reflect different amenities that these firms might offer to workers. 
\end{itemize}
\end{frame}
\begin{frame}
\label{cchk}
\frametitle{Card, Cardoso, Heining and Kline (2018)}
\begin{itemize}
\item Wages are set by employers to maximize profits, subject to labour supply constraints. 
\medskip
\item No search frictions, No ``unemployment".
 \medskip
 \item Workers value employers differently. This gives firms ``market power".
  \medskip
  \item CCHK model provides an important theoretical framework to link together rent sharing, monopsony models and firm inequality literature. 
\end{itemize}
\end{frame}
%%
\begin{frame}
\frametitle{Setup}
\begin{itemize}
\item $J$ firms and two types of workers: $L$ and $H$. 
\medskip
\item Each firm \underline{posts} a pair of wages $(w_{L_{j}},w_{H_{j}})$
 \medskip
 \item Workers observe the wages (costlessly). 
  \medskip
  \item Firms provide different ``workplace" environments over which workers have heterogenous preferences
  \be
  u_{i,S,j} = \beta_{S}\log(w_{S_{j}}-b_{S})+a_{S_{j}} + \epsilon_{i,S,j}
  \ee
  \item   $u_{i,S,j}$ indirect utility of worker $i$ in skill group $S$ if she works for firm $j$.
    \item   $w_{S_{j}}$ wage offered by firm $j$ to individuals in group $S$.
        \item   $b_{S}$ outside offer (e.g. wage paid in competitive industry).
         \item   $\epsilon_{i,S,j}$ idiosyncratic preferences (Type-I distributed). 
\end{itemize}
\end{frame}
%%
\begin{frame}[shrink=20]
\frametitle{Setup}
\begin{itemize}
\item Probability that a worker of type $S$ will show up at firm $j$ if this firm sets wage $w_{S_{j}}$?
\be
p_{S_{j}} \equiv \Pr\left(\argmax_{k\in\{0,\hdots,J\}}u_{iSk}=j\right) =\dfrac{\exp(\beta_{S}\log(w_{S_{j}}-b_{S})+a_{S_{j}})}{\sum_{k=1}^{J}\exp(\beta_{S}\log(w_{S_{k}}-b_{S})+a_{S_{k}})}
\ee
\item Trick: If $J$ is large then logit probabilities are well approximate by exponential probabilities
\be
p_{S_{j}} = \lambda_{S}\exp(\beta_{S}\log(w_{S_{j}}-b_{S})+a_{S_{j}})
\ee
where $\lambda_{S}$ constant common across all firms.
\medskip
\item If there is a total of $\mathcal{L}$ of low-skilled workers, then the labor supply of low-skilled workers that firm $j$ faces is given by
\be
\log(L_{j}(w_{L_{j}}))=\log(\lambda_{L}\mathcal{L}) + \beta_{L}\log(w_{L_{j}}-b_{L}) + a_{L_{j}}
\ee
Similarly for high-skilled workers
\be
\log(H_{j}(w_{H_{j}}))=\log(\lambda_{L}\mathcal{H}) + \beta_{H}\log(w_{H_{j}}-b_{H}) + a_{H_{j}}
\ee
\end{itemize}
\end{frame}
%%%%
%%
\begin{frame}
\frametitle{Firm Optimization}
\begin{itemize}
\item Firms cannot observe workers' preferences but have knowledge of their supply function.
\medskip
\item Assume that firms are price takers in the product market for the moment
\medskip
\item Output is produced according to
\be
Y_{j} = T_{j}f(L_{j},H_{j})
\ee
\medskip
\item where f(.,.) has CRS and we let $f_{S}\equiv\dfrac{\partial f(L_{j},H_{j})}{\partial S}$ for $S\in\{H,L\}$
\medskip
\pause
\item Firm problem
\be
\min_{w_{L_{j}},w_{H_{j}}} w_{L_{j}}L_{j}(w_{L_{j}}) + w_{H_{j}}H_{j}(w_{H_{j}}) \qquad s.t. T_{j}f(L_{j}(w_{L_{j}}),H_{j}(w_{H_{j}})) \geq \bar{Y}
\ee
\end{itemize}
\end{frame}
%%%%
\begin{frame}
\frametitle{Back to something familiar...}
\begin{itemize}
\item Taking FOC and rearranging, we obtain that
\be
w_{S_{j}} = \dfrac{1}{1+\beta_{S}}b_{S} + \dfrac{\beta_{S}}{1+\beta_{S}}T_{j}f_{S}\underbrace{\mu_{j}}_{\text{Marginal Cost of production}}
\ee
\medskip
\item What happens as $\beta_{S}\rightarrow\infty?$
\medskip
\item Who is capturing rents here and why? 
\medskip
\item Are amenities influencing the wage? Is this a standard compensating differential story?
\end{itemize}
\end{frame}
%%%%
\begin{frame}
\frametitle{A simple benchmark: linear production function}
\begin{itemize}
\item Suppose that we have a linear production function and that a firm faces a fixed output price $P_{j}^{0}$.
\be
Y_{j} = T_{j} N_{j}=T_{j}((1-\theta)L_{j}+\theta H_{j})
\ee
\medskip
\item Expression for wage of low skilled then becomes:
\be
w_{L_{j}} = \dfrac{1}{1+\beta_{L}}b_{L} + \dfrac{\beta_{L}}{1+\beta_{L}}T_{j}P_{j}^{0}(1-\theta)
\ee
\item Further assume that $b_{L}=(1-\theta)b$ and $b_{H}=\theta b$. Then
\be
\ln w_{L_{j}} = \ln \dfrac{(1-\theta)b}{1+\beta_{L}} + \ln(1+\beta_{L}R_{j})
\ee
\be
\ln w_{H_{j}} = \ln \dfrac{(1-\theta)b}{1+\beta_{H}} + \ln(1+\beta_{H}R_{j})
\ee
where $R_{j}\equiv \dfrac{T_{j}P_{j}^{0}}{b}$
\end{itemize}
\end{frame}
%%%%
\begin{frame}
\frametitle{Discussion}
\begin{enumerate}
\item No dynamics. 
\bigskip
\item No interactions $\rightarrow$ but if it is a small market a shock to a firm's productivity will affect all other firms! 
\bigskip
\item No compensating differentials! (will come back to this)
\end{enumerate}
\end{frame}
%%%%
\begin{frame}
\frametitle{What is the punch line?}
\begin{itemize}
\item Monopsonistic posting wage models $\rightarrow$ firms suddenly matters to explain wage inequality $\rightarrow$ typically neglected in SBTC and S-D-I models
\medskip
\item These ``firm effects" are going to be highly correlated with measures of firm productivity.
\medskip
\item Firm effects only arise due to the presence of ``frictions" in the labor market.  
\medskip
\item Many things can generate monopsony power of firms!
\medskip
\item Key insight: we can infer degree of market power of firms by looking at magnitude of pass-through of ``productivity" shocks into wages. 
\medskip
\item Example that do this really well:  Kline, Petkova, Williams, Zidar (2019); Garin and Silverio (2023)
\end{itemize}
\end{frame}
%%%%
\begin{frame}[shrink=10]
\frametitle{What if firm effects in pay reflect differences in work environments?}
\pause
\begin{itemize}
\item Utility from a job is not just a function of the wage but also of amenities offer by a given employer
\bigskip
\item In the standard CCHK model with linear production function amenities do not affect the marginal product of labor $\rightarrow$ no \underline{direct} effect on wages. Also, amenities are ``sunk" to the firm, i.e. the firm in CCHK model set only the wage not the level of amenity.
\bigskip
\item More realistic scenario is that firms might decide to offer a particular amenity (e.g. health insurance) and as result could pay a lower wage $\rightarrow$ compensating differential.
\bigskip
\item If that is the case, it would mean that firm effects in pay do not only reflect differences in productivity across firms but also differences in ``work environments" / compensating differential motives.
\bigskip
\item How much of the variability in firm-pay comes from unobserved amenities?
\bigskip
\item This is the research question of Sorkin (2018).
\end{itemize}
\end{frame}
\begin{frame}
\frametitle{But how to measure amenities?}
\begin{equation}
U_{ij} = \underbrace{\psi_{j} + \phi_{1}a_{j1}+\hdots+\phi_{K}a_{jK}}_{\text{Systematic Utility$\equiv V_{j}$}} + e_{ij} 
\end{equation}
\begin{itemize}
\item Hard to get data on all the amenities that influence the systematic utility of being employed by firm $j$. 
\medskip
\item Sorkin (2018) develops a way to infer $V_{j}$ without having data on the amenities (nor $\psi_{j}$) $\rightarrow$ application of the PageRank of algorithm used by Google to rank webpages.
\medskip
\item Intuition: if a lot of workers \textit{voluntary} leave employer $j$ to join employer $j'$, then it must be that employer $j'$ gives higher systematic utility to workers. 
\end{itemize}
\end{frame}
\begin{frame}[shrink=20]
\frametitle{Google Page Rank Applied to Job-To-Job transitions}
\begin{itemize}
\item Suppose you have a worker deciding on whether to work for either firm $j$ or firm $k$.
\medskip
\item Let $V_{j}$ represent the systematic value of being employed by firm $j$. Similarly for $V_{k}$.
\medskip
\item Worker's decision is based on comparing $V_{j}$ and $V_{k}$ as well a match-specific idiosyncratic factor assumed to be Type-1 distributed.
\medskip
\item The probability that the worker chooses $k$ over $j$ is thus given by
\begin{equation}
\dfrac{\exp(V_{k})}{\exp(V_{k})+\exp(V_{j})}
\end{equation}
\item In the data we observe $M_{kj}:$ number of people that left firm $j$ to join firm $k$. Similarly, we have $M_{jk}$, number of people that left firm $k$ to join firm $j$. 
\medskip
\item Then, letting $N$ represent the total number of workers, we can estimate the \textit{relative} common value of firm $k$ relative to $j$ as
\begin{equation}
\dfrac{M_{kj}}{N}\dfrac{N}{M_{jk}}=\dfrac{\exp(V_{k})}{\exp(V_{j})}
\end{equation}
\medskip
\item Can you see some problems with this?
\end{itemize}
\end{frame}
\begin{frame}[shrink=30]
\frametitle{Google Page Rank Applied to Job-To-Job transitions}
\begin{itemize}
\item Note that we can rewrite previous equation as
\begin{equation}
M_{kj}\exp(V_{j})=M_{jk}\exp(V_{k})
\end{equation}
\item This holds for all employers $j$ that we observe in the data. Let $J$ denote the set of firms
\pause
\medskip
\item We can now sum RHS and LHS across all employers
\begin{equation}
\underbrace{\sum_{j\in J}M_{kj}}_{\text{Number of workers leaving a given firm $j\in J$ to join $k$}}\underbrace{\exp(V_{j})}_{\text{Value of origin firm $j$}}=\underbrace{\sum_{j\in J}M_{jk}}_{\text{Number of workers leaving firm $k$}}\underbrace{\exp(V_{k})}_{\text{Value of firm $k$}}
\end{equation}
\item Therefore the value of firm $k$ is given by the following recursive expression
\begin{equation}
\exp(V_{k}) = \dfrac{\sum_{j\in J}M_{kj}\exp(V_{j})}{\sum_{j\in J}M_{jk}}
\end{equation}
\item In matrix form:
\begin{equation}
\exp(V) = S^{-1}M\exp(V)
\end{equation}
where $M$ is a $J\times J$ matrix with element $(j,k)$ given by $M_{jk}$ and $S$ is a diagonal matrix with kth entry given by ${\sum_{j\in J}M_{jk}}$; finally $\exp(V)$ is a $J\times 1$ vector that stacks together all the common utility values of the $J$ firms. 
\pause
\medskip
\item How we call that last equation? 
\medskip
\item This is what Google does when ranking webpages!
\end{itemize}
\end{frame}
\begin{frame}
\frametitle{37\% of all voluntary job transitions imply a decrease in pay!}
\begin{columns}
	\column{\dimexpr\paperwidth+5pt}
	\begin{center}
		\includegraphics[scale=0.35]{aux/sorkin}
	\end{center}
\end{columns}
\end{frame}
\begin{frame}
\frametitle{But 66\% of all transitions lead to a higher $V_{j}$!}
\begin{columns}
	\column{\dimexpr\paperwidth+5pt}
	\begin{center}
		\includegraphics[scale=0.35]{aux/sorkin2}
	\end{center}
\end{columns}
\end{frame}
\begin{frame}
\frametitle{Working in Mining: "standard" utility but need to pay higher than average to compensate for bad working conditions}
\begin{columns}
	\column{\dimexpr\paperwidth+5pt}
	\begin{center}
		\includegraphics[scale=0.35]{aux/sorkin3}
	\end{center}
\end{columns}
\end{frame}
\begin{frame}[shrink=10]
\frametitle{Summing up}
\begin{itemize}
\item Sorkin uses a stylized model where it shows that compensating differential motives account for around half the variability in firm-pay.
\medskip
\item How he arrived at that estimate from his model is what you guys have to solve in Ps-4!
\medskip
\item Punch-line: firm might offer different levels of pay \underline{and} utility.
\medskip 
\item This creates ``vertical differentiation" across employers: good firms might pay higher wages and offer better amenities while some firms offer bad pay and bad working conditions. 
\medskip
\item ``Frictions" permit the co-existence of these two types of firms in the labor market.
\medskip
\item But what are these frictions exactly? Is it a problem of asymmetric information? (See Jaeger et al. paper on how well workers know their outside options)
\medskip
\item Which factors create variability in firm-utility across employers offering the same level of pay? (Lachowska, Mas, Saggio, Woodbury: look at hours worked; Lachowska, Sorkin, Woodbury: Unemployment Insurance; Dube-Naidu-Reich: dignity of the workplace). 
\end{itemize}
\end{frame}
%%%%
\section{Applications}
\begin{frame}
\frametitle{Card, Cardoso and Kline (2016)}
\begin{itemize}
\item What explains the gender wage gap?
\medskip
\item Traditional view: women less skilled than men.
\begin{itemize}
\item But women have higher levels of education!
\end{itemize}
\medskip
\item Alternative view: women don't ask (Babcock and Laschever, 2009).
\pause
\medskip 
\item CCK: role of firms in explaining the gender wage gap?
\medskip
\item Two channels:
\medskip
\begin{enumerate}
\item Women sort in firms with lower premiums.
\item Women bargain a smaller wage effect premium than men. 
\end{enumerate} 
\end{itemize}
\end{frame}
%%
\begin{frame}
\frametitle{Portugal and the Gender Wage Gap}
\begin{columns}
	\column{\dimexpr\paperwidth+5pt}
	\begin{center}
		\includegraphics[scale=0.4]{aux/cck1}
	\end{center}
\end{columns}
\end{frame}
%%
\begin{frame}
\frametitle{Women have more schooling!}
\begin{columns}
	\column{\dimexpr\paperwidth+5pt}
	\begin{center}
		\includegraphics[scale=0.4]{aux/cck2}
	\end{center}
\end{columns}
\end{frame}
%%
\begin{frame}
\frametitle{CHK Event studies Men/Women in Portugal}
\begin{columns}
	\column{\dimexpr\paperwidth+5pt}
	\begin{center}
		\includegraphics[scale=0.4]{aux/cck3}
	\end{center}
\end{columns}
\end{frame}
%%
\begin{frame}
\frametitle{Wage change for men 1-4 Quartile of co-workers wages}
\begin{columns}
	\column{\dimexpr\paperwidth+5pt}
	\begin{center}
		\includegraphics[scale=0.4]{aux/cck4}
	\end{center}
\end{columns}
\end{frame}
%%
\begin{frame}
\frametitle{Wage change for women 1-4 Quartile of co-workers wages}
\begin{columns}
	\column{\dimexpr\paperwidth+5pt}
	\begin{center}
		\includegraphics[scale=0.4]{aux/cck6}
	\end{center}
\end{columns}
\end{frame}
\begin{frame}
\frametitle{Same gains?}
\begin{columns}
	\column{\dimexpr\paperwidth+5pt}
	\begin{center}
		\includegraphics[scale=0.4]{aux/cck8}
	\end{center}
\end{columns}
\end{frame}
%%
\begin{frame}{Oaxaca-Blinder Decompositions }

\begin{itemize}
\item {\small{}Standard tool in labor economics: the Oaxaca (1973)-Blinder
(1973) decomposition\medskip{}
}{\small\par}
\item {\small{}Let's briefly review how OB decompositions work\medskip{}
}{\small\par}
\item {\small{}There are two groups, $M$'s and $W$'s\medskip{}
}{\small\par}
\item {\small{}Average outcomes are $\bar{Y}_{M}$ and $\bar{Y}_{W}$, where
$\bar{Y}_{g}=E\left[Y_{ig}|group_{i}=g\right]$\medskip{}
}{\small\par}
\item {\small{}We hope to explain group differences with an observed covariate
vector $X_{i}$ }{\small\par}
\end{itemize}
\end{frame}

\begin{frame}{Oaxaca-Blinder Decompositions }

\begin{itemize}
\item {\small{}Quantity to be explained:}{\small\par}
\end{itemize}
\begin{center}
{\small{}$\Delta\equiv\bar{Y}_{M}-\bar{Y}_{W}$}{\small\par}
\par\end{center}
\begin{itemize}
\item {\small{}Suppose we run a regression for each group}{\small\par}
\end{itemize}
\begin{center}
{\small{}$Y_{ig}=X_{ig}^{\prime}\beta_{g}+\epsilon_{ig}$}{\small\par}
\par\end{center}
\begin{itemize}
\item {\small{}$X_{ig}$ includes a constant\medskip{}
}{\small\par}
\item {\small{}OLS coefficient vector:}{\small\par}
\end{itemize}
\begin{center}
{\small{}$\beta_{g}=E\left[X_{ig}X_{ig}^{\prime}|group_{i}=g\right]^{-1}E\left[X_{ig}Y_{ig}|group_{i}=g\right]$}{\small\par}
\par\end{center}

\end{frame}

\begin{frame}{Oaxaca-Blinder Decompositions }

\begin{itemize}
\item {\footnotesize{}By construction OLS fits group means:}{\footnotesize\par}
\end{itemize}
\begin{center}
{\footnotesize{}$\bar{Y}_{g}=\bar{X}_{g}^{\prime}\beta_{g}$.}{\footnotesize\par}
\par\end{center}
\begin{itemize}
\item {\footnotesize{}Therefore we can write}{\footnotesize\par}
\end{itemize}
\begin{center}
{\footnotesize{}$\Delta=\bar{X}_{M}^{\prime}\beta_{M}-\bar{X}_{W}^{\prime}\beta_{W}$}{\footnotesize\par}
\par\end{center}

{\footnotesize{}}

\begin{center}
{\footnotesize{}$=\left(\bar{X}_{M}-\bar{X}_{W}\right)^{\prime}\beta_{M}+\bar{X}_{W}^{\prime}\left(\beta_{M}-\beta_{W}\right)$}{\footnotesize\par}
\par\end{center}
\begin{itemize}
\item {\footnotesize{}The OB decomposition splits $\Delta$ into a component
explained by $X$'s, and a component explained by $\beta$'s\medskip{}
}{\footnotesize\par}
\item {\footnotesize{}First term answers the question: How much more would
M's make than W's if both groups were paid like M's for observables?\medskip{}
}{\footnotesize\par}
\item {\footnotesize{}Second term answers the question: How much of $\Delta$
is due to differences in earnings for M's and W's with the same characteristics?}{\footnotesize\par}
\end{itemize}
\end{frame}

\begin{frame}{Oaxaca-Blinder Decompositions }

\begin{center}
{\small{}$\Delta=\left(\bar{X}_{M}-\bar{X}_{W}\right)^{\prime}\beta_{M}+\bar{X}_{W}^{\prime}\left(\beta_{M}-\beta_{W}\right)$}{\small\par}
\par\end{center}
\begin{itemize}
\item {\small{}Can also write the alternative decomposition:}{\small\par}
\end{itemize}
\begin{center}
{\small{}$\Delta=\left(\bar{X}_{M}-\bar{X}_{W}\right)^{\prime}\beta_{W}+\bar{X}_{M}^{\prime}\left(\beta_{M}-\beta_{W}\right)$}{\small\par}
\par\end{center}
\begin{itemize}
\item {\small{}New first term answers the question: How much more would
M's make than W's if both groups were paid like W's for observables?\medskip{}
}{\small\par}
\item {\small{}Second term measures much of $\Delta$ is explained by the
$\beta$'s, weighting with M's characteristics}{\small\par}
\end{itemize}
\end{frame}
%%
\begin{frame}
\frametitle{AKM meets Oaxaca-Blinder}
\begin{itemize}
\item Gender specific firm effects model:
\begin{equation}
y_{it} = \alpha_{i} + \psi^{S}_{J(i,t)} + W_{it}^{\top}\gamma + r_{it} 
\end{equation}
\medskip
\item $\psi^{S}_{j}$ firm effect of firm $j$ for workers with sex =$\{M,F\}$
\medskip
\item CCK fit AKM models separately for men and women.
\item Dual Connected set: connected set of firms that overlap in both male and female AKM model.
\end{itemize}
\end{frame}

\begin{frame}
\frametitle{AKM meets Oaxaca-Blinder}
\begin{itemize}
\item Decompose differences in firm effects as
\begin{equation}
\begin{aligned}
& \E[ \psi^{M}_{J(i,t)} \vert Male] - \E[ \psi^{F}_{J(i,t)} \vert Female] = \\
& = \underbrace{\E[ \psi^{M}_{J(i,t)} - \psi^{F}_{J(i,t)} \vert Female]}_{\text{Bargaining Channel}} + \\
 & + \underbrace{\E[ \psi^{M}_{J(i,t)} \vert Male] - \E[ \psi^{M}_{J(i,t)} \vert Female]}_{\text{Sorting Channel}}
\end{aligned}
\end{equation}

\item First piece: difference in firm effects b/w men and women averaged using the distribution of jobs held by women.
\bigskip
\item Second piece: difference in average $\psi^{M}_{J(i,t)}$ across jobs held by women vs. men.
\end{itemize}
\end{frame}
%%
\begin{frame}
\frametitle{Alternative decomposition}
\begin{itemize}
\item Decompose differences in firm effects as
\begin{equation}
\begin{aligned}
& \E[ \psi^{M}_{J(i,t)} \vert Male] - \E[ \psi^{M}_{J(i,t)} \vert Female] = \\
& = \underbrace{\E[ \psi^{M}_{J(i,t)} - \psi^{F}_{J(i,t)} \vert Male]}_{\text{Bargaining Channel}} + \\
 & + \underbrace{\E[ \psi^{F}_{J(i,t)} \vert Male] - \E[ \psi^{F}_{J(i,t)} \vert Female]}_{\text{Sorting Channel}}
\end{aligned}
\end{equation}

\item First piece: difference in firm effects b/w men and women averaged using the distribution of jobs held by men.
\bigskip
\item Second piece: difference in average $\psi^{F}_{J(i,t)}$ across jobs held by women vs. men.
\end{itemize}
\end{frame}
%%
\begin{frame}
\frametitle{Gender Specific AKMs}
\begin{columns}
	\column{\dimexpr\paperwidth+5pt}
	\begin{center}
		\includegraphics[scale=0.35]{aux/cck9}
	\end{center}
\end{columns}
\end{frame}
%%
%%
%%
\begin{frame}
\frametitle{Gender Firm Effects and Productivity}
\begin{columns}
	\column{\dimexpr\paperwidth+5pt}
	\begin{center}
		\includegraphics[scale=0.45]{aux/cck11}
	\end{center}
\end{columns}
\end{frame}
\begin{frame}
\frametitle{Gender Firm Effects and Productivity}
\begin{columns}
	\column{\dimexpr\paperwidth+5pt}
	\begin{center}
		\includegraphics[scale=0.45]{aux/cck12}
	\end{center}
\end{columns}
\end{frame}
\begin{frame}
\frametitle{Gender Firm Effects and Productivity}
\begin{columns}
	\column{\dimexpr\paperwidth+5pt}
	\begin{center}
		\includegraphics[scale=0.45]{aux/cck13}
	\end{center}
\end{columns}
\end{frame}

%%
%%
\begin{frame}
\frametitle{Decomposition}
\begin{columns}
	\column{\dimexpr\paperwidth+5pt}
	\begin{center}
		\includegraphics[scale=0.4]{aux/cck14}
	\end{center}
\end{columns}
\end{frame}
\begin{frame}
\frametitle{Heterogeneity}
\begin{columns}
	\column{\dimexpr\paperwidth+5pt}
	\begin{center}
		\includegraphics[scale=0.34]{aux/cck15}
	\end{center}
\end{columns}
\end{frame}
\begin{frame}
\frametitle{Life Cycle}
\begin{columns}
	\column{\dimexpr\paperwidth+5pt}
	\begin{center}
		\includegraphics[scale=0.2]{aux/sorting_time}
	\end{center}
\end{columns}
\end{frame}
%%
\begin{frame}
\frametitle{Summing up}
\begin{itemize}
\medskip
\item Male/Female firm effects are highly correlated.
\medskip
\item But women only obtain 90\% of the firm effect of men.
\medskip
\item Firm effects explain 20\% of raw gender gap.
\medskip
\item Mostly explained by sorting component: women work in firms with lower premiums.
\medskip
\item But bargaining important in explaining the gap at higher levels of education.
\end{itemize}
\end{frame}
%%
\begin{frame}
\frametitle{Goldschmidt and Schmeider (2017)}
\begin{itemize}
\item Weil (2014): 
\begin{quote}
Wage discrimination is rarely seen in large  firms despite
the benefits it could confer. As long as workers are under
one roof, the problems presented by horizontal and vertical
equity remain. But what if the large employer could wage
discriminate by changing the boundary of the  firm?
\end{quote}
\medskip
\item By how can firms change their boundaries? Solution: by outsourcing. 
\medskip
\item GS studies on-site outsourcing events in Germany. Occupations studied: food, cleaning, logistics and services.
\medskip
\item Outsourcing Event$\equiv$ a large group of workers leave a mother firm and joins a new daugther firm in the business service.  
\end{itemize}
\end{frame}
%%
\begin{frame}
\frametitle{Temp agencies and Business Service Firms on the rise}
\begin{columns}
	\column{\dimexpr\paperwidth+5pt}
	\begin{center}
		\includegraphics[scale=0.2]{aux/gs_1}
	\end{center}
\end{columns}
\end{frame}
%%
\begin{frame}
\frametitle{AKM for FCLS and non-FCLS highly correlated}
\begin{columns}
	\column{\dimexpr\paperwidth+5pt}
	\begin{center}
		\includegraphics[scale=0.3]{aux/gs_2}
	\end{center}
\end{columns}
\end{frame}
%%
\begin{frame}
\frametitle{Large Wage Losses following the Outsourcing Event}
\begin{columns}
	\column{\dimexpr\paperwidth+5pt}
	\begin{center}
		\includegraphics[scale=0.3]{aux/gs_3}
	\end{center}
\end{columns}
\end{frame}
%%
%%
\begin{frame}
\frametitle{Losses almost all explained by lost of AKM effect}
\begin{columns}
	\column{\dimexpr\paperwidth+5pt}
	\begin{center}
		\includegraphics[scale=0.2]{aux/gs_4}
	\end{center}
\end{columns}
\end{frame}
\begin{frame}
\frametitle{Daruich, Di Addario, Saggio (2023): Temp Workers}
\end{frame}
\begin{frame}
\centering
\includegraphics[scale=0.2]{aux/ryan.jpg}
\end{frame}
\begin{frame}
\centering
\includegraphics[scale=0.2]{aux/ryan1.jpg}
\end{frame}
\begin{frame}
\frametitle{Large Wage Returns from Conversion}
\centering
\includegraphics[scale=0.2]{aux/JMP.jpg}
\end{frame}
\begin{frame}
\frametitle{Almost Entirely Explained by a \textit{Change} in Rent-Sharing Coefficients}
\centering
\includegraphics[scale=0.2]{aux/JMP1.jpg}
\end{frame}
%%
\begin{frame}
\frametitle{Drenik, Jaeger, Plotkin and Schoefer: Temp Workers in Argentina}
\centering
\includegraphics[scale=0.4]{aux/pascuel.jpg}
\end{frame}
%%
\begin{frame}
\frametitle{Taking Stock}
\begin{itemize}
\item \textit{Where} you work seems crucial to explain many import effects that we see in the labor market:
\begin{itemize}
    \item Losses following outsourcing.
    \item Gains from Temporary to Permanent conversion.
\end{itemize}
\pause
\item Let's now apply this logic to a topic covered in the previous set of slides: wage losses from job-displacement. 
\end{itemize}
\end{frame}

%%
\begin{frame}
\frametitle{Washington State}
\begin{columns}
	\column{\dimexpr\paperwidth+5pt}
	\begin{center}
		\includegraphics[scale=0.2]{aux/lmw}
	\end{center}
\end{columns}
\end{frame}
%%
%%
\begin{frame}
\frametitle{Germany}
\begin{columns}
	\column{\dimexpr\paperwidth+5pt}
	\begin{center}
		\includegraphics[scale=0.2]{aux/till_new}
	\end{center}
\end{columns}
\end{frame}
\begin{frame}
\frametitle{Evidence from multiple countries: Berthau, Acabbi, Barcelo, Gulyas, Lombardi, Saggio (forthcoming) }
\begin{columns}
	\column{\dimexpr\paperwidth+5pt}
	\begin{center}
		\includegraphics[scale=0.2]{aux/saggio}
	\end{center}
\end{columns}
\end{frame}

%%
\begin{frame}[shrink=7]
\frametitle{Summing up}
\begin{itemize}
\item Increasing availability of matched employer-employee datasets has helped economists assessing importance of firms in wage settings.
\medskip
\item More and more papers on this:
\medskip
\item Importance of firms in explaining wage ``gaps"
\medskip
\begin{itemize}
\item Male vs. Female.{\tiny{Card, Cardoso and Kline (2016)}}
\item College vs. Non-College. \\ {\tiny{Engdom, Moser (2017)}; Huneeus, Neilson, Miller and Zimmerman (2015)}
\item Black vs. White. {\tiny{Gerard, Lagos, Severnini, Card (2018)}} 
\item Layoff workers vs Non-Layoff workers. \\ {\tiny{Lachowska, Mas, Woodbury (2017)}; Schmieder, von Wachter and Heining (2018)}
\item Outsourced vs. Non Outsourced workers. {\tiny{ Goldschmidt and Schmieder (2017)}} 
\item Temporary vs. Permanent contract workers. {\tiny{Daruich, Di Addario, Saggio (2020)}; Drenik, Jager, Plotkin, Schoefer (2020)}
\item Unions vs. Non-Unionized Jobs  {\tiny{Beauregard-Lemieux-Messacar-Saggio (2024)}}
\item Family vs. non-family firms.   {\tiny{Di Porto-Pagano-Pezone Schivardi-Saggio (2024)}}
\end{itemize}
\item Lots of paper also on the microfoundations on how firms can impact wages beyond what shown in this slide-deck:  
{\tiny{Lamadon, Mogstad and Setzler (2020); Haanwinckel (2023); Berger Herkenhoff and Mongey (2019); Di Addario-Kline-Saggio-S\o lvsten (2023)}}.
\end{itemize} 
\end{frame}
%%
\begin{frame}
\begin{columns}
	\column{\dimexpr\paperwidth+5pt}
	\begin{center}
		\includegraphics[scale=0.6]{aux/meme}
	\end{center}
\end{columns}
\end{frame}
%%

\end{document}
